%% kapa: exam=1 round=25 year=2014 period=2 form=A subject=accounting subject_ko=회계학 seq=12 number=52
%% source: Q-Net 감정평가사 기출문제 (cst00309) artlSeq=5101207 file=extracted/1008/5101207/2014년 제25회 감정평가사 1차시험 기출문제/2교시(A형).hwp
%% conversion: kordoc 4.24.51 -> tools/gen1.py (auto); review=proofread-sonnet; flags=none
\begin{question}{52}
(주)한영의 20×1년의 기초재고액은 ￦200,000, 당기매입액은 ￦500,000이며, 기말 실사 결과 창고에 보유 중인 재고자산은 ￦50,000이다. 당기 말 현재 재고자산 관련 추가 자료는 다음과 같다. \par
{\small\centering
\begin{tabular}{|c|c|c|}
\hline
\makecell{항목} & \makecell{원가} & \makecell{비고} \\
\hline
\makecell{미착상품} & \makecell{￦30,000} & \makecell{선적지 인도조건으로 현재 운송 중이다.} \\
\hline
\makecell{적송품} & \makecell{100,000} & \makecell{해당 금액 중 수탁자가 40\%를 판매완료한 상태이다.} \\
\hline
\makecell{시용품} & \makecell{30,000} & \makecell{해당 금액 중 고객이 매입의사표시를 한 금액은 ￦10,000 이다.} \\
\hline
\makecell{재구매조건부판매} & \makecell{40,000} &  \\
\hline
\makecell{저당상품} & \makecell{20,000} & \makecell{기말재고자산 실사에 포함되었다.} \\
\hline
\makecell{반품가능판매} & \makecell{60,000} & \makecell{반품액의 합리적 추정이 불가능하다.} \\
\hline
\end{tabular}\par}

위 자료를 이용할 때 매출원가는? \par
\begin{choices}
  \item ￦400,000
  \item ￦420,000
  \item ￦440,000
  \item ￦460,000
  \item ￦500,000
\end{choices}
\end{question}
\answer{3}
